\documentclass[12pt,a4paper,openany]{book}

\usepackage{tocloft}
\usepackage{fontspec}
\setmainfont{Libertinus Serif}
\usepackage[greek.ancient,spanish]{babel}
\babelfont[greek]{rm}{Libertinus Serif}
\usepackage[a4paper,margin=3.2cm]{geometry}
\usepackage{microtype}
\usepackage[hidelinks]{hyperref}

% Griego politónico escrito directamente en UTF-8:
% \gr{...} para palabras sueltas; bloques con \foreignlanguage{greek}{...}
\newcommand{\gr}[1]{\foreignlanguage{greek}{#1}}
\newcommand{\stanzabreak}{\par\vspace{0.6\baselineskip}}

\setlength{\parindent}{1.25em}
\setlength{\parskip}{0.35em}
\raggedbottom


\hypersetup{
  pdftitle={Leer la Odisea en tiempos iletrados},
  pdfauthor={Juan José Gómez Cadenas}
}

\title{Leer la \emph{Odisea} en tiempos iletrados}
\author{Juan José Gómez Cadenas}
\date{}

\newcommand{\odiseachapter}[2]{\chapter{#2}}
\renewcommand{\thechapter}{\Roman{chapter}}
\setlength{\cftchapnumwidth}{3em}

\begin{document}

\frontmatter
\maketitle
\tableofcontents

\mainmatter

\odiseachapter{Introducción}{De libros y cine}

En julio de 1982, yo acababa de terminar el cuarto curso de la licenciatura en Ciencias Físicas y había ganado una beca para hacer prácticas de verano en la central nuclear de Cofrentes, por entonces todavía en construcción. Las prácticas empezaban en unos pocos días y aquella mañana regresaba de las oficinas de Hidroeléctrica (hoy Iberdrola) en Valencia, donde había firmado el primer contrato de mi vida laboral. De camino a mi piso de estudiantes, me entretuve en una de las librerías universitarias cercanas a Blasco Ibáñez, con el ánimo de hacerme con una buena novela para entretenerme durante los viajes entre Requena (donde estaba previsto que me alojara) y Cofrentes. Entre ambos pueblos había una hora de carretera. No teníamos redes sociales ni teléfonos móviles. Yo leía mucho, es cierto, pero leer no era nada tan extraño en los ochenta. El autobús que tomaría cada día durante los siguientes dos meses iba lleno de trabajadores de todos los oficios: electricistas, mecánicos, oficinistas, ingenieros y estudiantes en prácticas como yo. Y casi todos llevábamos un libro a cuestas para el trayecto.

El que casi me asaltó desde el escaparate de la librería aquella mañana fue nada menos que \emph{El señor de los anillos}. Eran tres volúmenes, editados por Minotauro, si la memoria aún sirve, y no eran baratos, pero los compré sin vacilar a cuenta de mi primer salario. Había oído hablar de Tolkien, porque alguien había contado una historia de las revueltas estudiantiles del mayo del 68 en París, cuando en el metro aparecían pintadas que afirmaban en rojo vivo sobre las paredes «Frodo vive». Yo no sabía muy bien quién era Frodo ni qué tenía que ver con la revolución, pero cualquier sombra de duda se despejó apenas leí los versos con los que arrancaba la trilogía.

\begin{verse}
Tres anillos para los reyes elfos bajo el cielo.\\
Siete para los señores enanos en palacios de piedra.\\
Nueve para los hombres mortales condenados a morir.\\
Uno para el señor oscuro, sobre el trono oscuro\\
en la tierra de Mordor donde se extienden las sombras.
\end{verse}

¿Por qué aquellos versos fueron suficientes para vaciarme los bolsillos y comprometerme con mil páginas de relato épico? Creo que la respuesta tiene que ver con el hecho de que mi cabeza estaba llena de palabras, leídas en las innumerables novelas y libros de poesía que ya había devorado por entonces. Palabras que resonaron con aquellas otras palabras del poema inicial, como un diapasón resuena al sonido de otro.

Mi primera gran novela, si se me permite la licencia poética, fue la Biblia, en concreto el Antiguo Testamento, en una versión ilustrada para niños que leía en casa de mi tía Corro con seis, quizás siete, años de edad. Era uno de esos libros de antes, de los que llevaban en su catálogo los vendedores ambulantes de enciclopedias, primorosamente encuadernado, con hojas de papel satinado y láminas en color cada pocas páginas. Contaba la historia de un pueblo, protegido por un Dios poderoso, algo arbitrario y no poco cruel, y los episodios eran pura acción y maravilla. Mares que se abrían al paso de las tribus y volvían a cerrarse sobre los ejércitos del faraón, murallas que se derrumbaban al resonar las trompetas, un profeta que hablaba con una zarza ardiente, una divinidad capaz de todo tipo de acciones fabulosas e incomprensibles: pedir a su devoto servidor que matara a su hijo y pararle la mano en el último momento, convertir a una mujer en una estatua de sal, enviar a una banda de ángeles asesinos a que mataran a los primogénitos de los enemigos de su pueblo.

Sin darme cuenta, en casa de mi tía Corro, estaba leyendo mi primera gran épica. No fue la única. Cada domingo, mi padre me compraba un tebeo (las aventuras del Capitán Trueno) y un libro de la colección «Historias Selección», donde descubrí a Robinson Crusoe, a Ben-Hur, a Tom Sawyer, a Nemo, a Uncas y a John Silver, entre otros muchos. A los quince años, cuando mi padre me regaló la \emph{Ilíada} y la \emph{Odisea} ---de esa historia hablaré en el próximo capítulo---, yo era ya una máquina de leer. A los veintidós, un libro de trescientas páginas caía en una semana, dos a lo sumo. Los tres volúmenes de \emph{El señor de los anillos} no fueron ninguna excepción.

La trilogía de Peter Jackson llegó veinte años más tarde. Naturalmente, hice cola en el cine para verla y, naturalmente, la he vuelto a ver muchas veces desde entonces. La vi cuando mi hija Irene era un bebé y mi hijo Héctor aún no había nacido, la vi otra vez cuando aún eran tan niños que los abuelos nos reñían por darles sustos con los Nazgûl, y otra vez cuando eran un poco más mayores, y otra vez de adolescentes y otra vez durante la covid y otra vez recientemente. Por culpa del pesado de su padre, mis hijos se saben la trilogía de memoria.

Y sin embargo, ninguno de los dos ha leído los libros.

Los libros han estado siempre en casa: la versión española que compré en 1982 y la versión en inglés que compré en una librería de viejo en San Francisco hacia 1989, durante mi posdoctorado en California. Están muy visibles en las estanterías, provocadores y sensuales como las sirenas, con su excelente encuadernación y sus mil páginas de prodigios. En la misma estantería están \emph{Las mil y una noches}, los poemas completos de Czesław Miłosz y los de Raymond Carver, \emph{El Aleph} de Borges, \emph{Rayuela} de Cortázar, \emph{Ana Karenina} de Tolstoi y \emph{Los desposeídos} de Ursula K. Le Guin. Todos esos volúmenes están esperando la misma atención que esperaban los libros de la biblioteca de mi padre. Pero yo, a aquellos, no los decepcioné. Los de la mía esperan en vano a Godot.

Irene es médica; Héctor, un brillante estudiante de físicas, acabando su máster. Los dos leían de pequeños. Son de la generación de Harry Potter y además en casa no teníamos televisión, así que, durante su niñez, como para su padre antes que ellos, los libros eran un excelente antídoto contra el aburrimiento. Ambos siguen leyendo hoy en día. Irene tiene poco tiempo, pero lo compensa con un extraño sexto sentido para seleccionar novelas interesantes. De hecho, fue ella la que me descubrió a una de mis escritoras favoritas, Elena Ferrante, y leímos juntos, gozosamente, la tetralogía de \emph{La amiga estupenda}. Héctor acaba de leer \emph{La insoportable levedad del ser} y los dos hemos pasado muchas horas hablando de la honda, casi inaudita percepción del alma humana que allí ofrece Milan Kundera.

¿Entonces? La diferencia entre mis hijos y yo no es de actitud, ni de interés. La diferencia está en los números. Aunque a ellos las redes sociales les sorprendieron relativamente tarde, no han sido inmunes a un asalto que ha mermado el tiempo de todo el mundo, sobre todo de la gente joven, robándoselo a muchas otras actividades más enriquecedoras para el alma, en particular la lectura. En el mundo en el que yo crecí, la televisión era una receta fija que ocupaba unas pocas horas del día, entre una carta de ajuste que abría a las siete o las ocho de la tarde y cerraba dos horas después de la cena. No había internet y el cine era un lujo reservado a los domingos. No había, por supuesto, TikTok ni Instagram. Yo leía antes de dormir, leía los fines de semana, en vacaciones, durante los viajes, a cualquier hora. No lo hacía por devoción intelectual o por ganar cultura. Lo hacía, a menudo, porque no había nada mejor que hacer. Pero una infancia y una adolescencia de lecturas a destajo construyen un vocabulario y una imaginería a la que uno ya no puede sustraerse.

Los elfos de \emph{El señor de los anillos}, los enanos y los orcos, los Rohirrim, las temibles batallas del Abismo de Helm y el viaje final de Frodo por Mordor se formaron en mi mente antes de ver una sola imagen en la pantalla. Para mis hijos y para casi todos los jóvenes del planeta, todos esos seres fantásticos, todas esas tierras misteriosas son las que Peter Jackson imaginó. En tiempos de libros, había tantos elfos como lectores. Hoy en día, Galadriel tiene el rostro de Cate Blanchett, Arwen es Liv Tyler, Frodo es Elijah Wood. Nada que objetar, todos ellos son magníficas encarnaciones de los personajes de la leyenda. Pero cuando el verbo se hace carne, irremediablemente, pierde parte de su carácter divino.

Mis hijos, como muchos jóvenes en la franja entre los veinte y los treinta años, han leído menos que sus padres, pero la lectura ha estado y sigue estando presente en sus vidas. Ahora bien, entre los de menor edad, es frecuente leer ---odiosos deberes aparte--- solo los \emph{feeds} de sus redes sociales. Aún más preocupante es un secreto a voces que la capacidad de comprender un texto complejo se ha desplomado entre los jóvenes. No menos sabido es el hecho de que su vocabulario se ha reducido. Y si combinamos ambas cosas, el resultado inevitable es una generación con menos recursos expresivos, menos capacidad analítica, menos instrumentos de raciocinio. O, en otras palabras, una generación de jóvenes con menos capacidad para pensar por sí mismos, por tanto, quizás más crédulos, menos críticos, más fáciles de manipular, más susceptibles a las estafas del populismo, las teorías de la conspiración o la demoledora tontería de los \emph{influencers}. La lectura, me temo, no es un lujo de señoritos, sino, como decía Gabriel Celaya de la poesía, un arma cargada de futuro. Los versos de Miguel Hernández anunciaban:

\begin{verse}
Tristes armas\\
si no son las palabras.\\
Tristes. Tristes.
\end{verse}

Y en tiempos iletrados, como los que vivimos, nos han robado las palabras.

Ahora bien, como afirma Rick Blaine ---tan indistinguible de Humphrey Bogart como Aragorn de Viggo Mortensen---, siempre nos quedará París, es decir, el cine. Cierto, mis hijos no han leído \emph{El señor de los anillos}, pero podría argumentarse que las películas son tan buenas que una cosa compensa la otra. Y quizás lo mismo ocurre con todas las grandes historias. No hay gran libro, se diría, que no cuente con su versión cinematográfica y, en algunos casos, el celuloide supera de largo a la página escrita. \emph{Blade Runner}, por ejemplo, le da vuelta y media a la novela original (\emph{Do Androids Dream of Electric Sheep?}) y la trilogía de \emph{El padrino} reside en una esfera del Parnaso a la que no habrían llegado jamás las novelas de Mario Puzo. \emph{Muerte en Venecia} (Mann) o \emph{Los muertos} (Joyce) son obras maestras en el papel y en el cine. La serie televisiva sobre \emph{La amiga estupenda} de Ferrante es tan buena como el texto escrito, entre otras cosas porque la autora participó activamente en el guion. La lista es larga y habría mucho que escribir sobre la intertextualidad del cine y la novela, sobre la miríada de formas en que se complementan y se completan.

Entonces, ¿podemos prescindir alegremente del anticuado papel que nos obliga a descifrar signos y construir imágenes en el aire? ¿Podemos fiarle nuestra cultura, es decir, nuestra formación como ciudadanos, al cine, a las series televisivas? ¿Se pierde algo esencial por no haber leído a Tolkien? Quizás los obsesos de la lectura son \emph{boomers} nostálgicos que no se enteran de que los tiempos han cambiado, como esos programadores chapados a la antigua que aún no han admitido que todo su conocimiento de C++, o de cualquier otro lenguaje de programación, es inútil en los tiempos de la IA.

¡Ay!, mucho me temo que cuando el cine ---o la serie televisiva--- reemplaza por completo a la lectura, hay algo que se queja. Cine y literatura tienen su espacio propio, sus recursos únicos ---también sus limitaciones---, y en mi opinión, forman una pareja perfecta, uno de esos matrimonios enamorados que son como un magnífico animal bicéfalo, cada parte de la pareja un individuo único, pero el total mayor que la suma de sus partes. Pero hay una asimetría. Durante tres mil años, los humanos han usado los libros para construir el mundo, mientras que el cine es un recién llegado, un chaval de poco más de un siglo. Se diría, sin embargo, que nos movemos hacia un mundo iletrado, un mundo donde no ha hecho falta quemar los libros como en \emph{Fahrenheit 451} porque nos hemos limitado a ignorarlos, un mundo donde toda la información y todo el entretenimiento van a ser audiovisuales, un mundo, de hecho, donde incluso el cine se va a ver amenazado. ¿Qué joven tiene paciencia para sentarse quieto dos o tres horas delante de una pantalla, salvo para seguir las aventuras de un adolescente en pijama (me refiero a Spider-Man, naturalmente)?

%La literaturaEl cine, inevitablemente, con pocas excepciones, desarrolla argumentos lineales y el espectador, a menudo, es dócilmente pastoreado a lo largo del planteo, nudo y desenlace. Se pierde, en consecuencia, uno de los mayores placeres de la lectura que podríamos denominar «avanzada» y que permite el lector reconstruir o descubrir la historia ---léase Rayuela de Cortazar, para un clásico, o la Historia del amor, de Nicole Krauss, para un ejemplo reciente---.

%De hecho, el cine genera la obsesión por evitar el malhadado \emph{spoiler}. No queremos que nos cuenten qué pasa en la peli (no digamos ya el final) porque el argumento se convierte en el motor principal de la historia. A ese respecto, confieso sin pudor que mis amigos y familiares me temen, por mi tendencia arruinarles el espectáculo. No puedo evitarlo, en general, el argumento de un relato me trae sin cuidado ---es una razón por la que nunca me han enganchado las novelas de misterio--- y de hecho, mi técnica para comprar libros de los que sé poco o nada, suele ser leer la primera página \emph{y la última}. Me importa como un autor acaba su relato, igual que me importa como lo empieza, porque los primeros y últimos párrafos dicen mucho de su estilo y de sus recursos literarios, pero me da igual enterarme de que Gatsby muere antes de empezar la novela, porque todo el valor literario de esa obra maestra no tiene nada que ver con el detalle, obvio por otra parte desde el primer renglón, de que el protagonista es un héroe trágico, con un final inevitable.

¿Qué nos da la literatura que es difícil de captar en el cine? Quizás el ritmo interno de la prosa ---Fitzgerald---, el estilo ---Borges, Nabokov---, la extrañeza ---Márquez, Cortázar---, la locura ---Bulgakov---. La sensación de que Flaubert ha sacado un conejo de la chistera cuando, para demostrar que el marido de Emma Bovary es un pobre hombre, no habla de él, sino de su sombrero. La intuición de que la realidad es un tenue velo que esconde algo más transcendente, que alimenta la lectura de Shakespeare. El simple y puro placer de repetir en voz alta unos versos perfectos.

\begin{verse}
Te acordaste de mí, cuando subías,\\
al silencio que sufre la serpiente\\
prisionera de grillos y de umbrías.
\end{verse}

¿Cómo reducir a celuloide el infinito que despliega Lorca en tres líneas? ¿Cómo expresar la emoción con que el corazón tiembla al leer «el silencio que sufre la serpiente»?

%Hay más. Trescientas páginas ---no digamos ya mil--- dan de sí para crear personajes muy complejos, cuya invocación requiere espacio ---que le pregunten, si no, a Proust---. Dan de sí para inventar todas esas pequeñas ramas que surgen del tronco del argumento principal y lo revisten de significado. Dan de sí para jugar con el tiempo, acelerándolo o deteniéndolo, como hace Fitzgerald en el primer capítulo de \emph{Gatsby}:
%
%\begin{quote}
%We walked through a high hallway into a bright rosy-colored space, fragilely bound into the house by French windows at either end. The windows were ajar and gleaming white against the fresh grass outside that seemed to grow a little way into the house. A breeze blew through the room, blew curtains in at one end and out the other like pale flags, twisting them up toward the frosted wedding cake of the ceiling---and then rippled over the wine-colored rug, making a shadow on it as wind does on the sea.
%
%The only completely stationary object in the room was an enormous couch on which two young women were buoyed up as though upon an anchored balloon. They were both in white and their dresses were rippling and fluttering as if they had just been blown back in after a short flight around the house. I must have stood for a few moments listening to the whip and snap of the curtains and the groan of a picture on the wall. Then there was a boom as Tom Buchanan shut the rear windows and the caught wind died out about the room and the curtains and the rugs and the two young women ballooned slowly to the floor\footnote{He preferido copiar el original, cuya belleza es apabullante. Añado aquí una traducción al español: \itshape Pasando por un corredor de techo alto llegamos a un alegre espacio de colores vivos, apenas integrado a la casa por ventanales franceses a lado y lado. Los ventanales blancos estaban abiertos del todo y resplandecían contra el césped verde de la parte de afuera, que parecía entrar un poco en la casa. La brisa soplaba a través del cuarto, haciendo elevarse hacia adentro la cortina de un lado y hacia afuera la del otro, como pálidas banderas, enroscándolas y lanzándolas hacia la escarchada
%cubierta de bizcocho de novia que era el techo, para después hacer rizos sobre el tapiz de color vino tinto, formando una sombra sobre él, como el viento al soplar sobre el mar.
%
%El único objeto completamente estacionario en el cuarto era un enorme sofá en el que había dos mujeres flotando como sobre un globo anclado. Ambas vestían de blanco, y sus trajes revoloteaban ondulados como si hubieran acabado de regresar por el aire tras un corto vuelo por los alrededores de la casa. Debí haber permanecido unos instantes escuchando el restañar y revolotear de las cortinas y el crujir del retrato de la pared. Entonces se escuchó una explosión cuando Tom Buchanan cerró el ventanal de atrás y el viento atrapado murió en el cuarto, y las cortinas y los tapetes y las dos mujeres descendieron cual globos con lentitud hasta el piso.}.
%\end{quote}

%Hay al menos dos excelentes versiones cinematográficas de \emph{Gatsby}, interpretadas por actores tan monumentales como Robert Redford (dirigido por Jack Clayton en 1974, sobre guion de Coppola) y Leonardo DiCaprio (en 2013). Ninguna de las dos consigue captar, ni de lejos, ese momento mágico, esa metáfora que asimila a Daisy y a Jordan Baker con ángeles o hadas, la misma que sugiere ---solo con el acto de cerrar los ventanales--- que Tom Buchanan es un personaje brutal.

%Igualmente difícil, aunque no imposible es captar las sutiles ironía que una novela puede entretejer. Un ejemplo que aún recuerdo de mis lecturas durante los años de la universidad fue el viaje siete, en «Diario de las estrellas, Viajes y Memorias», de Stanislav Lem. El protagonista, Ijon Tichy, viaja solo en su nave espacial, en la vecindad de la estrella
%Betelgeuse cuando un meteorito perfora el casco de su nave, destrozando el timón, que tiene que reparar para salir del apuro. Pero se da la circunstancia que hacen falta al menos dos personas para ese arreglo, una para apretar tornillos y otra para sostener las tuercas. Por fortuna, la nave a la deriva entra en una zona donde los vértices gravitatorios crean distorsiones del espacio tiempo cuyo resultado es que la nave se llena de dopplegangers de Tichy. Nada más fácil que ponerse de acuerdo, ahora que son multitud para arreglar el timón, pero en lugar de hacerlo, las copias de Tichy se enzarzan en una discusión estéril que abarca desde la naturaleza del libro albedrío hasta quién se come el último pedazo de chocolate. En el fondo, ninguno de los dopplegangers quiere ser el que sostenga la tuerca... hasta que dos versiones infantiles del astronauta resuelven el problema, mientras los adultos siguen empecinados en su inútil asamblea.
%
%¿Se puede hacer una versión cinematográfica de esta historia? Sin duda, pero el resultado ---que yo sepa no se ha intentado--- será posiblemente una historia más gráfica ---sin duda igual de divertida, con la nave llena de copias del astronauta hablando todos a la vez--- y menos filosófica. Los cuentos de Lem, desternillantes, invitan además a pensar. El riesgo de la pantalla es quedarnos con el chiste y olvidar la reflexión.

Puede argumentarse, y con razón, que el cine tiene sus propias sutilidades, su propia magia ---esa combinación de teatro, música y diálogo que permite escenas sin equivalente en la página escrita---, su propia capacidad de emocionarnos y sobrecogernos, gracias a las maravillosas bandas sonoras, a las cámaras audaces, a los efectos especiales. Puede argumentarse, y con razón, que el cine se merece de sobra el calificativo de séptimo arte\footnote{De hecho, mi adicción al cine es casi tan temprana como mi vicio por la lectura. A los diez años, mi familia se mudó a un pueblo de la costa mediterránea, situado entre Vinaroz y Amposta, al que siempre me he referido como Alba, quizás porque, como Cervantes, prefiero no acordarme de su nombre. Mi padre era la autoridad militar de la plaza, lo que nos daba a sus hijos la identidad de extraterrestres y a mí me granjeó durante mucho tiempo el estatus de paria. Durante la semana, me veía a menudo con Quelo, mi único amigo, pero los domingos Quelo tenía otros planes, con sus amigos de siempre. Yo solo tenía la sesión doble en el cine del pueblo y una bolsa de pipas, para olvidar la soledad.}.
% y aseguró e Pero para mí, el riesgo más grande del cine no está en su capacidad de emular o no, a la literatura que traduce. En mi opinión, ambos se complementan y se completan. He nombrado antes La amiga estupenda y el ejemplo no puede venir más a cuento. La serie televisiva, en cuatro temporadas, es una lectura a la vez fiel y libre de las novelas. No es sorprendente, ya la autora, Elena Ferrante, fue una de las guionistas (junto con Francesco Piccolo, Laura Paolucci y Saverio Costanzo), pero sí muy gratificante. Hay muchos ejemplos similares que podrían darnos cierta tranquilidad. El cine se tiene ganado el título de séptimo arte, después de todo.

Pero el cine es también un negocio que mueve muchísimo dinero. Una película taquillera, como la versión de Spider-Man del verano de 2026, ha recaudado más de dos mil millones de dólares. Por comparación, los cuatro grandes detectores del Gran Colisionador de Hadrones (LHC) del CERN cuestan unos 1.500 millones de dólares; se pagaron en un periodo de diez años mediante el esfuerzo colectivo de una colaboración internacional que incluye miles de científicos y centenares de laboratorios, y están entre los aparatos más caros y sofisticados que ha producido la ciencia.

Y una película taquillera tiene que gustarle a millones de personas, lo que implica, además de actores famosos y mucho espectáculo, un argumento que no ponga nervioso a nadie. Por ejemplo, un jovencito con poderes arácnidos, dando fabulosos saltos y combatiendo a malos suficientemente domesticados.

¿Pero qué ocurre cuando un ambicioso director se propone narrar una de las grandes épicas de la literatura? Sabemos, por \emph{El señor de los anillos}, que es posible hacerlo, pero conviene caer en la cuenta de que Tolkien, grandioso como es, no es muy complicado. Los buenos son buenísimos y carecen de fisuras; los malos comen niños fritos y son feísimos; el talante moral de Aragorn, de Arwen, de Gandalf o de Legolas, es el mismo que el del Capitán Trueno; Galadriel ofrece una pizquita de oscuridad semejante a la aceituna en el martini de bondad de los protas. Frodo y Sam protagonizan un \emph{bromance} impecable; Boromir se hace el malo dos minutos y a continuación muere por la causa. El único personaje complejo es Gollum, que, por otra parte, acaba por redimirse y sacar a Frodo de apuros. Con toda mi admiración hacia Jackson, la trilogía de Tolkien se dejaba querer.

%¿Y la Ilíada? La película Troya (2004, Wolfgang Petersen), es una lectura libre, comercial, inexacta e inteligente de la gran épica, apta para todos los públicos, con las dosis justas de gore y un Brad Pitt que clava el papel de Aquiles. Tiene mucho de espectáculo y poca profundidad, pero aún así consigue rescatar algunas de las escenas más relevantes del gran poema homérico, incluyendo la batalla de los héroes ---narrada solo a medias, como veremos--- y la generosidad final de Aquiles con Príamo. Pero los personajes, por supuesto, están fabricados con alambre y las deformaciones narrativas llegan a ser tan grotescas como inventarse la muerte de Agamenón y Menelao, un sacrilegio equivalente a matar a Frodo, o a D'artagnan. No contento con ello, el travestido director indulta a Paris y manda a Helena camino de Roma. Troya, claramente, no es una adaptación de la Ilíada, pero tampoco es una relectura. No añade ningún valor a la historia ni a los personajes y sustrae mucho. El espectador que ha leído a Homero pasa el mismo buen rato que viendo un Western. El joven que no lo conoce, se va a su casa con un catálogo de falsedades, una historia dulzona y unos personajes más planos que el cartón.

¿Y la \emph{Odisea}? En este caso ha habido varios intentos con diferentes méritos y limitaciones. El \emph{Ulises} de Camerini, en 1954, protagonizado por Kirk Douglas, tenía escenas afortunadas ---incluyendo la del cíclope--- y trucos muy innovadores, como usar la misma actriz para encarnar a Penélope y Circe. Recuerdo haber visto esta película en el cine de Alba, con doce o trece años y haberme pasado muchas tardes jugando a ser Ulises y Telémaco, a tensar el arco y asaetear a los pretendientes, empezando por Antínoo, que se parecía mucho al Churrero, el abusón que me traía por la calle de la amargura en el instituto. Naturalmente, era demasiado joven para entretenerme en sutiles análisis de contenido, pero agradezco de todo corazón la compañía de Kirk Douglas en mis juegos solitarios.

En 1968, Franco Rossi se propuso rodar una versión de la \emph{Odisea} a la altura de la épica de Homero. Empezó por escoger el formato apropiado, esto es, el de una miniserie, con un total de ocho capítulos, cada uno de una hora. Rossi entendió perfectamente que la extensión era imprescindible si se pretendía un mínimo de fidelidad a una historia densísima, en la que cada episodio, cada personaje, cada encuentro contribuyen a la construcción de la épica. A continuación, los guionistas adaptaron con bastante fidelidad el texto de Homero, incluyendo la mayoría de los episodios relevantes de la épica: las andanzas completas de Telémaco, Polifemo, sirenas, Circe, Calipso, Nausícaa, la corte de los feacios, el viaje a los infiernos y, por supuesto, la matanza de los pretendientes.

La serie tiene sesenta años, así que puede argumentarse que viene a cuento una nueva versión, una nueva lectura del mito que aproveche los muchos recursos tecnológicos de los que disponemos, como el célebre IMAX. Claro que también sería deseable que la voluntad pedagógica que mostraron Rossi y la RAI se perpetúe, sobre todo si tenemos en cuenta los argumentos sobre los que venimos reflexionando desde hace un rato. En nuestra época iletrada, el único contacto con el mito de Odiseo que van a tener millones de personas es el que se derive de esa nueva serie o película, tal como ya sucedió, hace más de dos décadas, con \emph{El señor de los anillos}.

Así que cuando el celebrado director Christopher Nolan anunció sus planes de rodar una nueva versión de la \emph{Odisea}, pensé que se trataba de buenas noticias. Con Nolan me pasa como con Joaquín Sabina. A veces me carga con su histrionismo, o me pone de mala leche con su tendencia a exhibirse. A veces me empacha tanto ego o me impaciento con sus manías. Pero no me pierdo una película suya, como nunca dejo de escuchar a Sabina. Da igual que el taimado cantautor sea a menudo un cínico con pedales, o un malote de barrio. Cuando quiere, y quiere muy a menudo, es capaz de codearse con Quevedo y con Lorca, con Cohen. Da igual que el bienfamoso director sea un místico \emph{new age}. Cuando quiere, también le salen escenas como la del teseracto en \emph{Interstellar}, que quitan el hipo.

El problema, como veremos, es que Homero es mucho más difícil de manejar que Tolkien. Odiseo es un personaje casi inabarcable, contradictorio y no poco amoral. Astuto, sagaz, mentiroso, artero, cruel, implacable, odioso, valiente, elegante, generoso a veces, buen hijo y mejor padre, mal marido, gran amante, observador, sensible, capaz de hacerse atar a un mástil para escuchar a las sirenas y de arrojar a un bebé desde las murallas de Troya.

Un personaje así presenta un problema serio al cine, sobre todo si se trata de un director \emph{superstar} que quiere rodar una película taquillera. Pero quizás, me dije a mí mismo, sir Christopher demostraría estar a la altura de tamaño reto y nos ofrecería, si no una lectura literal ---imposible en tres horas de película--- al menos una interesante versión del héroe que enriqueciera el mito.

La película se estrenó en julio de 2026, a la vez que salía a imprenta la maravillosa y atrevida traducción de la \emph{Odisea} de Juan Manuel Tobal, publicada por Hipálage. Confieso que la lectura de esa traducción fue, literalmente, un empujón al túnel del tiempo. Me vi otra vez con quince años, leyendo por primera vez a Homero, declamando las líneas de Segalá a grito limpio, persiguiendo a mi primo Pedro con el libraco abierto que me había regalado mi padre, leyendo como quien lee la Biblia. Cuando conseguí volver al presente, me di cuenta de que Tobal había conseguido exactamente lo que se proponía al atreverse con una nueva traducción.

¿Y qué había conseguido?

Hacer magia.

Unos pocos días después, fui a ver la \emph{Odisea} de Nolan. A lo largo del último mes, la he visto tres veces. He leído las mil y una críticas que se le han dedicado. Los elogios me han parecido forzados en la mayor parte de los casos, aunque no en todos ---hay escenas, como el momento en que se oye la música de las sirenas, que van a pasar a la historia del cine---. Las críticas, a menudo, también. Me he asombrado de la tinta vertida en cuestiones irrelevantes, como el color de la piel de Helena ---cuando es evidente que Lupita Nyong'o es tan divinal como corresponde al personaje---, el realismo de las armaduras ---¿a quién le importa?--- o el hecho de que el caballo de Troya no tuviera ruedas ---Homero no dice que las tenga, por cierto---. Creo, en cambio, que no se ha hablado bastante de lo que la película de Nolan dice de los hipócritas, mojigatos y, me temo, iletrados tiempos que nos ha tocado vivir.

\emph{Leer la Odisea en tiempos iletrados} surgió como una reacción inmunitaria a la picadura del áspid y la del mosquito. Dejo al lector que decida cuál es cuál en las obras a las que me he referido.

Hay otra razón más que me ha llevado a escribir este ensayo. En 2024, publiqué mi propia lectura de la \emph{Odisea}, imaginando a un Ulises octogenario que repasa su vida en una Ítaca que ya no siente como suya. Es un anciano lúcido, cansado y carcomido por los remordimientos, pero a diferencia del personaje que nos dibuja Nolan, no sufre de estrés postraumático y no necesita drogas en forma de flores de loto ni psicoanalistas olímpicas. Por el contrario, es capaz de asumir el tipo de hombre que es y tomar una decisión inesperada, que lo libera de los fantasmas.

Dichas las motivaciones, el objetivo principal de este libro es invitar (y acompañar) al lector en la lectura de Homero, a través de la espléndida traducción de Tobal. En cuanto a mi admirado sir Christopher, uso y abuso de él ---bueno, de su peli--- sin pudor y con no poca maldad, por lo que pido de antemano perdón a todos los olímpicos, porque no creo que Nolan sea el problema. El problema, como veremos, es que el cine actual ---el cine comercial, destinado a entretener, que no a educar a las masas--- teme a Homero. El problema es que el Odiseo de Homero nos pone un espejo delante en el que se refleja nuestra propia crueldad y falta de escrúpulos, y por eso el director más popular del momento lo ha convertido en un cliché santurrón, reflejando nuestra infinita hipocresía.

\input{src/odisea1_nolan}
\input{src/odisea2_osada}
\input{src/odisea3_helena}
\odiseachapter{Menelao}{Menelao}\label{cap:menelao}

En el capítulo anterior, protestaba de la pobre lectura que nos ofrece Nolan de Helena. Eso no quiere decir ---he insistido ya en el tema en los capítulos precedentes--- que el afamado director estuviera obligado a ofrecer la misma versión de la divina que presenta Homero. De hecho, puede argumentarse que el único sentido que tiene revisitar las épicas clásicas es reinventárselas. Mi queja no es que Nolan reinterprete a Helena, sino que sale del paso con un cliché.

Permíteme, paciente lectora, generoso lector, que ilustre lo que quiero decir ---con toda desvergüenza--- usando un fragmento nuevo de \emph{Abandonando Ítaca}.

En la versión publicada del texto, Helena y Menelao aparecen solo de refilón. La escritura del capítulo anterior, sin embargo, me inspiró una extensión al poema, que ofrezco aquí a la generosa lectora, al paciente lector. Ulises empieza por recordar la (imaginaria) reunión anual, en Esparta, de los héroes que todavía quedan vivos, todos ellos ya muy ancianos.

\begin{versopropio}
Uno de los mejores momentos del año\\
es cuando viajas a Esparta.\\
Menelao organiza una reunión, a la que acuden\\
los que sobrevivieron a Troya y al \emph{nóstos}\footnote{El regreso a la patria después de la guerra.}.
\stanzabreak
Por supuesto, eres el invitado de honor,\\
te agasajan con comida y vino,\\
y hay una linda esclava de suaves manos\\
que cuida bien de tu cuerpo cansado.
\stanzabreak
Y todos disfrutáis, recordando\\
episodios repetidos tantas veces con los años\\
que es difícil separar la verdad de la fantasía,\\
la realidad de la imaginación.
\stanzabreak
A veces sospechas que la mitad de tus hazañas\\
no son más que reconstrucciones, tan lindas como tu esclava,\\
de algo más simple y sucio y brutal,\\
conjuradas para cubrir la sangre y la mugre y el estiércol.
\stanzabreak
¿Pero a quién le importa ya? Os reunís cada año,\\
y cada año hay menos y menos invitados en casa de Menelao,\\
y coméis, y bebéis, y cantáis, para celebrar\\
el fingimiento de que vuestras vidas tuvieron algún sentido.
\end{versopropio}

Tras lamentarse de que la mayoría de sus compañeros han muerto, Odiseo reflexiona sobre los dioses:

\begin{versopropio}
Cada vez sois menos. La mayoría ya no está.\\
Muchos de tus amigos no volvieron de Troya,\\
Aquiles, Áyax, Patroclo, Antíloco,\\
la lista es larga. Otros murieron en el viaje de vuelta.
\stanzabreak
Mucho se ha escrito sobre los dioses\\
oponiéndose al regreso de los argivos,\\
y poco sobre las ratas que, sospechas,\\
estaban detrás de las enfermedades que mataron a tantos.
\stanzabreak
Ah, pero a la gente le encanta cantar a los inmortales,\\
ruines y caprichosos como son,\\
y celebrar sus enredos en los asuntos humanos,\\
su lujuria y su crueldad dementes.
\stanzabreak
Puestos a elegir, tú te quedarías con las ratas,\\
que, al menos, iban a lo suyo,\\
y aunque se comían tu comida y se bebían tu agua,\\
no se divertían con tu ruina, como los olímpicos.
\end{versopropio}

Odiseo no solo está desengañado de los dioses. Con el tiempo, se ha vuelto ateo y su falta de fe le ha revelado también cuál es la verdadera utilidad de los olímpicos.

\begin{versopropio}
Hace ya mucho tiempo\\
que dejaste de creer en esa gente\\
supuestamente divina, que bebía néctar y daba fiestas\\
en lo alto del monte Olimpo.
\stanzabreak
Lo más probable es que ninguno de tus compañeros,\\
los pocos viejos supervivientes aún en pie,\\
creyera en ellos. Pero todos fingís\\
una devoción apropiada para mantener el statu quo.
\stanzabreak
Y es que, si los dioses no son reales,\\
¿por qué habría la gente de tomarse en serio a los reyes?\\
¿Por qué sufrir a tiranos como vosotros,\\
si están tan vacíos de sentido como las deidades?
\stanzabreak
Así que seguís con las quemas de muslos,\\
y las libaciones, y las plegarias, diligentemente conducidas\\
por esa otra especie de ratas, los sacerdotes,\\
siempre listos para comerse tu comida y beberse tu vino.
\stanzabreak
A cambio, por supuesto, de profecías\\
que resultan coincidir con tus propios planes,\\
y señales que confirman tus antojos,\\
y una bendición general de tus caprichos.
\stanzabreak
Así que, es verdad, los dioses no existen,\\
ni son necesarios, después de todo:\\
los hombres pueden encargarse de abusar,\\
y esclavizar, y matarse unos a otros.
\end{versopropio}

¿Quién, entre las lectoras de cierta edad, entre los lectores que ya peinan canas, no recuerda esos reencuentros de la promoción del instituto, o de la facultad, esas fiestas bienintencionadas para celebrar los veinticinco, o los treinta años del nosotros, los de entonces? ¿Y quién no ha caído en la cuenta de que nosotros, los de entonces, ya no somos los mismos? ¿Quién no se ha asustado con las arrugas de la más guapa de la clase, la barriga cervecera del galán de antaño, los ojos cansados y miopes del otrora poeta, la voz áspera del que cantaba como Orfeo? ¿Quién no ha reconocido su propia decadencia en la de aquellos amigos de otros tiempos?

\begin{versopropio}
Los que todavía se reúnen cada año\\
fingen ser hermanos de sangre,\\
y es cierto, hay calor entre vosotros,\\
algo parecido a la amistad.
\stanzabreak
Ah, pero con qué avidez\\
calibra cada invitado el estado del otro,\\
mide la piel arruinada,\\
el pelo menguante, la espalda rota.
\stanzabreak
Con qué precisión evaluáis\\
los ojos apagados, la mano temblorosa,\\
los párpados caídos, el temblor en la voz,\\
todas las señales que gritan vejez.
\stanzabreak
Este ha engordado demasiado, este está demasiado flaco,\\
Idomeneo se pasa el día durmiendo, Filoctetes siempre está borracho,\\
Néstor apenas puede hablar, la mayoría\\
está en peor forma que tú.
\stanzabreak
Y te preguntas si el único sentido de esas reuniones\\
es reinventar el pasado,\\
espantar la soledad que avanza,\\
y contar los cuerpos supervivientes.
\end{versopropio}

Finalmente, Odiseo reflexiona sobre su amigo Menelao:

\begin{versopropio}
Menelao se conserva bastante bien,\\
puede argumentarse que él y tú\\
sois los menos dañados por el tiempo,\\
sobre todo porque ambos seguís lúcidos.
\stanzabreak
Y él es un verdadero amigo,\\
hay algo de verdad entre las muchas mentiras,\\
fuisteis hermanos de armas,\\
y más de una vez os salvasteis la vida el uno al otro.
\stanzabreak
Así que el vínculo es real,\\
tan real, de hecho, que solo Telémaco\\
es más querido para tu corazón,\\
sin contar, por supuesto, a Circe.
\stanzabreak
Pero tú no le mencionarás la hechicera\\
a tu amigo, como él no mencionará a Helena,\\
los dos habéis aprendido\\
a respetar el corazón roto del otro.
\stanzabreak
No es que seáis adolescentes enamorados,\\
las penas de los jóvenes se arreglan tan fácilmente\\
como sus cuerpos fuertes. Pero el amor perdido de un viejo\\
es el último dolor que llevará a la tumba.
\end{versopropio}

En el siguiente fragmento, Helena ya no está presente en la corte del viejo Menelao:

\begin{versopropio}
Y así, Menelao no habla de Helena,\\
hace ya tantos años que se fue,\\
de hecho, desapareció poco después\\
de la visita de Telémaco a Esparta.
\stanzabreak
Tu amigo presumió tanto ante tu hijo,\\
se aseguró de que viera a la legendaria Helena en todo su esplendor,\\
para que se lo contara a su padre Odiseo,\\
por si acaso había sobrevivido.
\stanzabreak
Y en efecto, Telémaco dijo\\
que nunca había visto a nadie más hermoso,\\
ni más distante. Mirando su fría majestad,\\
se convenció de su naturaleza divina.
\stanzabreak
Así que cuando Helena se desvaneció, el relato oficial\\
habló de una santa ascensión al Elíseo,\\
no se había fugado otra vez, no,\\
sino que los dioses la habían reclamado.
\end{versopropio}

Los versos siguen la narrativa de la \emph{Odisea} y a la vez la subvierten. Siguen así:

\begin{versopropio}
Nadie cuestionó la historia,\\
como nadie cuestiona tu aventura con el cíclope,\\
los cuentos se inventaron precisamente para esto,\\
para encontrar sentido donde solo hay ruido.
\stanzabreak
Ruido y furia, como las hordas que condujiste,\\
rabia, y odio y lujuria, bombeando en vuestros corazones,\\
bronce, y sangre y fuego,\\
para ser disfrazados de batallas épicas.
\stanzabreak
Así que fuisteis todos a Troya,\\
a vengar el rapto de la esposa de Menelao,\\
y luchasteis diez años, y quemasteis la ciudad,\\
y devolvisteis el trofeo al palacio.
\stanzabreak
Pero la segunda vez que ella desapareció,\\
no había ciudad que saquear, ni unidad entre los argivos,\\
y demasiados viejos al mando,\\
así que invocar a los dioses resultó más apropiado.
\end{versopropio}

Odiseo reflexiona sobre el primer rescate de Helena, más motivado quizás por la codicia del botín que por la necesidad de recuperar a la adúltera, y muestra la futilidad de toda la empresa. La primera desaparición da la excusa para una guerra que, seguramente, habría ocurrido de todas maneras; cuando llega la segunda, no pasa nada.

\begin{versopropio}
Lo asombroso\\
es que Menelao la amaba.\\
La amó mucho más después de Troya\\
de lo que la había amado nunca antes.
\stanzabreak
Después de la boda,\\
Helena era solo un trofeo, y que el príncipe,\\
tan joven y furioso, se llevara a la cama a otras mujeres\\
era solo buen deporte, legal, para un hombre.
\stanzabreak
El supuesto secuestro lo enfureció,\\
quizás porque sabía, como todo el mundo,\\
que nadie había obligado a Helena a saltar\\
a la nave de Paris, ni a la cama de Paris.
\stanzabreak
Cuentan los relatos que en Troya\\
Helena desnudó el pecho ante su esposo,\\
lo invitó a golpear, y Menelao\\
dejó caer la espada y lloró.
\stanzabreak
Y por una vez estás convencido\\
de que esta historia sí es cierta:\\
Menelao había necesitado un rechazo y una guerra\\
para entender que estaba enamorado.
\end{versopropio}

El amor a Helena llega tarde y mal y es por ello mucho más trágico. La guerra deja de ser una excusa y se convierte en una justificación y más tarde en un absurdo.

\begin{versopropio}
Pero entonces cometió el error clásico:\\
asumió que su amor recién descubierto,\\
su deseo ferviente, su admiración por ella,\\
le ganarían el corazón.
\stanzabreak
Se equivocaba. Helena, por supuesto,\\
se comportaría como la noble dama que era,\\
y se ocuparía de la casa y del lecho,\\
guardándose sus cartas.
\stanzabreak
Pero fue su esposo quien mató a Deífobo,\\
y tú la viste llorar con el alma rota sobre su cuerpo,\\
qué ironía, pensaste. Ni el valiente Menelao,\\
ni el cobarde Paris, sino aquel otro,
\stanzabreak
príncipe discreto y valeroso,\\
casi tan bravo y honesto como Héctor.\\
A él, y solo a él,\\
le dio Helena su corazón.
\end{versopropio}

Helena, como cualquier personaje literario, admite infinitas lecturas. La mía aquí le hace elegir no al hermoso ---y cobarde--- Paris, ni al aguerrido ---e inútilmente enamorado--- Menelao, sino al más humilde de los héroes, el callado, gentil Deífobo. Toda historia de amor que se precie es una tragedia y la de Helena y Menelao no es una excepción. Mientras el príncipe aqueo bebe los vientos por su esposa, ella trama, fría y pacientemente, su venganza:

\begin{versopropio}
Y por él, y solo por él,\\
tramó su venganza.\\
No eligió el cuchillo, como su hermana,\\
ni el veneno que tan bien conocía.
\stanzabreak
No: ella quería herir más hondo,\\
así que trabajó duro hasta que Menelao\\
solo podía respirar el aire que ella había respirado,\\
y entonces lo apuñaló por la espalda.
\stanzabreak
No hizo falta bronce, ni siquiera las nuevas dagas de hierro.\\
Ella, simplemente, se fue;\\
dicen algunos que un troyano llamado Eneas\\
organizó la expedición de rescate.
\end{versopropio}

Así que en la vuelta final de tuerca, Helena huye con Eneas, rompiendo, esta vez sí, el corazón de Menelao:

\begin{versopropio}
Menelao no habla de ella,\\
como tú no hablas de Circe,\\
cada año, los dos os reunís\\
en Esparta y bebéis vino espeso,
\stanzabreak
y os emborracháis, y lloráis por el amigo perdido,\\
y por el amor perdido\\
que nunca merecisteis.
\end{versopropio}

A eso me refiero cuando hablo de reinventar un personaje como Helena. Nadie pide fidelidad a Homero, pero sí estar a la altura de sus personajes.

\input{src/odisea4_calipso}
\input{src/odisea5_nausicaa}
\input{src/odisea6_esqueria}
\input{src/odisea6_laguerra}
\input{src/odisea7_ciclope}
\input{src/odisea8_eolo}
\input{src/odisea8_circe}
\input{src/odisea9_hades}
\input{src/odisea10_agamenon}
\input{src/odisea11_sirenas}
\input{src/odisea12_pretendientes}
\input{src/odisea13_ninas}
\input{src/odisea14_penelope}
\input{src/odisea15_homero}
\input{src/odisea_marcharse_de_itaca}

\end{document}
